\subsection{对应的性质 V：Kac--Moody 表示}

\ssszh 我们已来到图 \ref{f:table} 的最后一行，这里会发生某种相当引人注目的事情。确切地说，我们将遇到对应的一条猜想性质，由 Frenkel--Gaitsgory \cite{frenkgaitslocallangds} 提出，它涉及 Kac--Moody 表示。这样的性质在 $p$-adic 情形中没有对应物，是局部朗兰兹中几何情形特有的一个侧面。

\begin{re} 我们应当提到，Kac--Moody 表示在几何朗兰兹中的重要性，最早出现在整体的（global）故事中。如我们在第 \ref{ss:hitchinhecke} 节所暗示的，Beilinson--Drinfeld 给出的一种构造自守 D-模的方法，就是通过 Kac--Moody 表示的局部化。这可以看作 Frenkel--Gaitsgory 工作的整体类比，也是「与共形场论的联系提供了算术情形所不具备的工具」这回事第一次出现。\end{re}

\ssszh 我们想先让读者对 Kac--Moody 代数有一个基本的感觉。对一般的连通约化群，它们的讨论本质上化归到 $\fg$ 单的情形。因此，为避免多半只是记号上的困扰，我们假定 $G$ 是单连通、几乎单的群。

首先从环路代数开始
\[
\fg(\!(t)\!) \simeq \fg \underset  k \otimes k(\!(t)\!).
\]
这不过是把 $\fg$ 作标量扩张到 Laurent 级数域。确切地说，括号的形式为
\[
[ X \otimes f, Y \otimes g] = [X,Y] \otimes fg, \quad X,Y \in \fg, \quad
f, g \in k(\!(t)\!).
\]
把它想成从穿孔圆盘 $\sD^\times$ 到 $\fg$ 的映射空间、而李括号逐点进行，会很有帮助。也就是说，它是穿孔圆盘上平凡 $G$-丛的无穷小对称。

我们把环路代数看作 $k$ 上的无穷维李代数。一个基本的事实是：许多乍看像是 $\fg(\!(t)\!)$-模的 $k$-向量空间，实际上带有中心扩张的作用，即仿射李代数的作用
\[
0 \rightarrow k \cdot \mathbf{1} \rightarrow \gk \rightarrow
\fg(\!(t)\!) \rightarrow 0.
\]

\begin{re} 这可以看作一般现象的一例：场论的量子化会给其局部对称性带来反常。\end{re}

幸运的是，环路代数的一切可能中心扩张都有简单的分类——它们与有限维代数 $\fg$ 上不变双线性形式所在的直线典范地一一对应。确切地说，对这样的形式 $\kappa \in (\fg^* \otimes \fg^*)^G$，相应地有中心扩张
\[
0 \rightarrow k \cdot \mathbf{1} \rightarrow \fgk \rightarrow \fg(\!(t)\!)
\rightarrow 0,
\]
其中由中心性，$\mathbf{1}$ 与任何元素的括号为零，其余括号则修正为
\[
[X \otimes f, Y \otimes g] = [X,Y] \otimes fg - \kappa(X,Y) \on{Res} (f
\cdot dg).
\]

\ssszh 对水平 $\kappa$，以 $\fgk\mod$ 记光滑 $\fgk$-模构成的范畴，其中中心元 $\mathbf{1}$ 通过恒等作用。参见例 \ref{ex:kmmods}。若 $\kappa$ 是整的，即环路代数的扩张来自 $LG$ 的中心扩张，那么如例 \ref{ex:kmmods} 所述，这个范畴带有 $\Dmod(LG)$ 的作用；因此下面只关注这一情形。

回忆 $G$ 被假定为单且单连通，整水平构成一个格
\[ \mathbb{Z} \hookrightarrow (\fg^*
\otimes \fg^*)^G,\]
由所谓的基本形式生成；具体说来，它使短余根的长度的平方为二。

\ssszh 仿射李代数表示论中的一个基本观察是：随着水平 $\kappa$ 变化，相应的表示论有两个「相」。

确切地说，对一切充分负定的 $\kappa$，表示论的表现类似。例如，所有 Verma 模都有有限长度，而这些水平上其单商的特征标公式本质上具有相同的形式 \cite{kashiwaratanisakinegativeintegrallevel}。

同样，对一切正定的 $\kappa$，表示论随水平变化的表现也类似。此时 Verma 模的长度无限，其单商的特征公式不随水平发生实质变化 \cite{kashiwara90}。

因此可以预期存在一个特殊水平 $\kappa_c$，即临界水平，表示论在此发生相变。朴素地看，人们或许以为这个中心点就是零形式。其实，由于 $\fg$ 表示论中 $\rho$-移位的仿射类比，它是 Killing 形式的一半再取负，即
\[
\kappa_c = - \frac{1}{2} \cdot \kappa_{\on{Killing}}.
\]

\ssszh 正如所料，相变处即临界水平处，表示论的行为显示出其他水平所没有的许多特征。

Feigin 与 Frenkel \cite{feiginfrenkel91} 发现了这样一条基本性质。非临界水平处，包络代数的中心\footnote{技术上说，我们指适当的完备包络代数的中心，或等价地，光滑 $\fgk$-模阿贝尔范畴的 Bernstein 中心。}是平凡的，即只有标量；而他们证明临界水平处的中心非平凡。关键在于，他们进一步把它等同于穿孔圆盘上 $\cG$-oper 的模空间 $\Op$ 上的函数代数。虽然我们马上会回忆它的精确定义，但首先的要点是：临界水平处的表示论把局部朗兰兹对偶、尤其是 $\cG$-联络的几何，硬性地编码于自身之中。

\ssszh 粗略地说，一个 $\cG$-oper 由一个带联络的 $\cG$-丛和一个 $\cB$-约化构成，二者适当地横截。

我们把它写清楚。回忆一下，为了谈论朗兰兹对偶群，需要固定 $G$ 与 $\cG$ 的带钉结构，特别地固定 $\check{\fg}$ 的 Chevalley 生成元 $f_i, \check{\alpha}_i, e_i, i \in I$。若以 $f$ 记 $\cfg$ 中使
\[
(f, 2 \check{\rho}, \underset i \Sigma \hspace{.2mm} e_i )
\]
构成 $\fsl_2$ 三元组的元素，则 $\cG$-oper 的空间是联络的模空间
\begin{equation} \label{e:defopers}
\Op \simeq \{ d  + f dt + \check{\fb}(\!(t)\!)dt \} / L\check{N}.
\end{equation}
把平凡丛看作带有其恒等的 $\cB$-约化，则上式中的 $f dt$ 一项就实现了上述联络与约化的横截性。关于 oper 定义的其他表述（例如具有明显不依赖坐标 $t$ 的优点），可参见 \cite{BD}。

公式 \eqref{e:defopers} 可以读作：取平凡 $\cG$-丛上全体联络 $d + \cfg(\!(t)\!)dt$ 的某个仿射子空间，再按 $L\check{N}$ 的伴随作用取商。事实上，\eqref{e:defopers} 是仿射空间 $d + \cfg(\!(t)\!)dt$ 关于 $L\check{N}$ 的一个通用特征标的 Hamilton 约化，特别地带有典范的 Poisson 结构。Gelfand--Dikii 与 Drinfeld--Sokolov \cite{drinfeldsokolov} 在几何朗兰兹出现之前，就在可积层级的背景下研究了它的 Poisson 几何；对 $\fg = \fsl_2$，这恢复了著名的 KdV 层级。

对尚未遇到过这个模空间的读者，这里有两条帮助定位的注。

\begin{re} 虽然 $\Op$ 与 $\LS$ 类似，被表示成联络的空间模去环路群的作用，但它的几何要简单得多，两者不可同日而语。事实上，它本身就是个无穷维仿射空间；特别地，它的点没有自同构。
\end{re}

\begin{re} 上一条注通常这样证明：说明每个 $L\check{N}$-轨道都是单传递的，即没有非平凡稳定化子，并给出与每个轨道恰好相交一次的显式截面。

指出整个这套说法（包括上述论证）是 Kostant \cite{kostantslice} 一项工作的仿射类比，会很有帮助。他证明了有典范同构
\[
\cfg /\!/ \cG \simeq (f + \check{\fb}) / \check{N},
\]
也就是说，对 $\cfg$ 作类似的有限维 Hamilton 约化，会切出 Poisson 中心的一个拷贝。此外，他证明两边的自然量子化是匹配的，即中心 $Z(\cfg)$ 可以由 $U(\cfg)$ 关于 $\cN$ 的量子 Hamilton 约化切出。\end{re}

\ssszh 通过投影映射 $\Op \rightarrow \LS$（从 oper 过渡到其底层的联络），可以把 $\QCoh(\Op)$ 看作 $\LS$ 上的范畴层。

于是人们预期，在局部几何朗兰兹下 $\gkc\mod$ 与 $\QCoh(\Op)$ 互换；从而恒等的幺半函子
\[
\QCoh(\Op) \rightarrow \Hom_{\TwoCoh(\LS)}(\QCoh(\Op), \QCoh(\Op))
\]
经由 Feigin--Frenkel 同构 $\OO_{\Op} \simeq Z(\gkc)$，与中心作用
\[
Z(\fgk)\mod \rightarrow \Hom_{\Dmod(LG)\mod}(\gkc\mod, \gkc\mod)
\]
互换。

在下面几个小节中，我们会讨论这个对应给出的、关于临界水平表示论的若干有趣猜想。

不过我们想在此先提一个特别引人注目的推论。已知任何 $\cG$-联络 $\sigma$ 都带有 oper 结构 $\chi$，即投影 $\Op \rightarrow \LS$ 在 $k$-点上满射 \cite{frenkelzhuoper}。由上述，形式上可推出：环群对应于 $\sigma$ 上摩天楼 $\Vect$ 的「不可约」范畴表示 $\sC_{\sigma}$，由
\begin{equation} \label{e:kmcenc}
\gk\mod \underset{\QCoh(\Op)} \otimes \Vect,
\end{equation}
给出，即中心特征标为 $\chi$ 的 Kac--Moody 表示范畴。

一方面，这通过 Kac--Moody 表示给出了全部 $\sC_{\sigma}$ 的统一猜想构造，即局部几何朗兰兹对应的谱到自守方向。另一方面，这预言：对任何两个位于 $\sigma$ 之上的 oper $\chi, \chi'$，相应的 Kac--Moody 模范畴等价——这本身就是表示论中一个深刻的猜想。

\begin{re} 事实上，由于尚未解释的、对朴素表述 \eqref{e:loclangnaive} 的修正，范畴 \eqref{e:kmcenc} 其实只是 $\sC_{\sigma}$ 的全子范畴。实践上，这意味着通过对 \eqref{e:kmcenc} 作重正规化来得到 $\sC_{\sigma}$，进一步的讨论参见第 \ref{ss:locaff} 节。\end{re}

\ssszh 这就结束了我们对图 \ref{f:table} 中出现的对象的引导性讨论。不过，在继续之前，我们应当强调：来自算术故事以及与物理的联系，还有更多预期的相容性，参见下文第 \ref{ss:variants} 节。

我们忍不住要提一个特别壮观的例子。Braverman--Finkelberg 对 $G = GL_n$ 提出了卷积代数 $\Dmod(LGL_n)$ 本身的朗兰兹对偶的一个候选 \cite{bravermanfinkelbergconjecturecoloumb}。这基于物理上的考虑，最终可追溯到 Gaiotto--Witten \cite{gaiottowitten} 发现的一个引人注目的公式，用于计算四维 $\mathscr{N}=4$ Yang--Mills 理论的 S-对偶边界条件。确切地说，他们猜想：在局部朗兰兹对偶下，$\Dmod(LGL_n)$ 与参数化局部系统退化旗的模空间上的凝聚层互换，即秩为 $i$ 的局部系统 $\sL_i$（$1 \leqslant i \leqslant n$）连同未必单射的联络映射
\[
\sL_i \rightarrow \sL_{i+1}, \quad 1 \leqslant i < n.
\]
此后 Zhu \cite{zhufamilies} 提出了这一猜想的算术类比。
